\documentclass[11pt]{article}

% ---------- Packages (per course conventions) ----------
\usepackage[margin=1in]{geometry}
\usepackage{graphicx}
\usepackage{float}
\usepackage{booktabs}
\usepackage{siunitx}
\usepackage{amsmath}
\usepackage{amssymb}
\usepackage{tikz}
\usepackage{tcolorbox}
\usepackage{enumitem}
\usepackage{xcolor}
\usepackage{hyperref}
\hypersetup{colorlinks=true, linkcolor=black, urlcolor=blue}

\setlength{\parindent}{0pt}
\setlength{\parskip}{\baselineskip}

% =========================================================
% EDIT THESE FOR EACH NEW ASSIGNMENT
% =========================================================
\newcommand{\assignmentNumber}{X}
\newcommand{\assignedDate}{}
\newcommand{\dueDate}{}
\newcommand{\numProblems}{8} % just for your own reference; add/remove \problempage calls below as needed

% ---------- Boxed-answer helper ----------
\newtcolorbox{finalanswer}{
    colback=white, colframe=black, boxrule=1pt,
    title=Answer, fonttitle=\bfseries, coltitle=black,
    colbacktitle=white
}

% ---------- Blank sketch box helper ----------
% Usage: \sketchbox[height]
\newcommand{\sketchbox}[1][6cm]{%
    \begin{tikzpicture}
        \draw[thick] (0,0) rectangle (\textwidth,#1);
        \node[anchor=north west, gray] at (0.1,#1-0.1) {\small \textit{Sketch / free-body diagram (label axes, forces, angles)}};
    \end{tikzpicture}
}

% ---------- Reusable problem scaffold ----------
% \problempage{Problem number}{Problem prompt/summary text -- paraphrase, don't copy verbatim}
\newcommand{\problempage}[2]{%
    \clearpage
    \section*{Problem #1}

    \textbf{Given / Find:} #2

    \vspace{0.3cm}
    \textbf{Sketch:}

    \sketchbox

    \vspace{0.5cm}
    \textbf{Equations \& Solution:}

    \vspace{0.2cm}
    % --- Set up governing equations here, e.g.: ---
    % \begin{align*}
    %     \sum F_x &= 0 \\
    %     \sum F_y &= 0
    % \end{align*}

    \vspace{4cm}

    \begin{finalanswer}
        % Report to 3 significant figures with units, e.g.: $F_R = \SI{393}{lb}$, $\theta = 353^\circ$
    \end{finalanswer}
}

\begin{document}

% =========================================================
% COVERSHEET
% =========================================================
\begin{titlepage}
    \centering
    \vspace*{2cm}
    {\LARGE \bfseries ENGR 225: Mechanics 1} \\[0.3cm]
    {\Large Homework Coversheet} \\[0.3cm]
    {\large Assoc.\ Prof.\ Clayton Byers} \\[2cm]

    \begin{flushleft}
        \Large
        \textbf{NAME:} \underline{\hspace{10cm}} \\[1cm]
        \textbf{DATE:} \underline{\hspace{10cm}} \\[1cm]
        \textbf{ASSIGNMENT:} \underline{\hspace{9cm}} \\[2cm]
    \end{flushleft}

    \large
    \textbf{Level of involvement of AI, LLMs, or related tools in the completion of this assignment:}

    \vspace{1cm}

    \begin{flushleft}
        \Large
        $\square$ \quad None \\[0.7cm]
        $\square$ \quad Minor \\[0.7cm]
        $\square$ \quad Major \\[0.7cm]
        $\square$ \quad Extensive \\[1.5cm]
    \end{flushleft}

    \normalsize
    \textit{Mark the box that applies (e.g., replace $\square$ with $\blacksquare$, or type an X beside it, before compiling/printing). Your disclosure does not affect your grade -- it keeps you accountable for the work you are doing.}

\end{titlepage}

% =========================================================
% HOMEWORK HEADER
% =========================================================
\section*{ENGR 225 -- Homework \assignmentNumber}
\textbf{Assigned:} \assignedDate \hfill \textbf{Due:} \dueDate

\vspace{0.3cm}
\textit{Formatting per course guidelines: each problem on its own page, with problem summary, labeled sketch, governing equations, units carried throughout, results to 3 significant figures, and a clearly boxed final answer.}

% =========================================================
% PROBLEMS
% Copy/paste the \problempage block below for each new problem,
% renumber the first argument, and fill in a short paraphrased
% Given/Find summary as the second argument. Delete unused blocks
% for assignments with fewer problems, or duplicate for more.
% =========================================================

\problempage{1}{
    % Paraphrase what's given and what you're solving for.
}

\problempage{2}{
    %
}

\problempage{3}{
    %
}

\problempage{4}{
    %
}

\problempage{5}{
    %
}

\problempage{6}{
    %
}

\problempage{7}{
    %
}

\problempage{8}{
    %
}

\end{document}